\documentclass{article}

\usepackage{float}
\restylefloat{table}

\usepackage{booktabs}

\title{Team Productivity: POC\\\progname}

\author{\authname}

\date{}

\input{../Comments.text}
\input{../Common.text}

\begin{document}

\maketitle

This document summarizes the contributions of each team member up to the POC
Demo.  The time period of interest is the time between the beginning of the term
and the POC demo.

\medskip

\noindent \textbf{Project URL:}
\url{https://github.com/AHMEDELZARIA/squigglr}

\noindent \textbf{Data collected:} November 19, 2026, from the issues, pull
requests and \texttt{main} branch of the project repository. The counts
exclude the issue and the pull request that deliver this report.

\section{Demo Plans}

\subsection*{Feasibility risk}

The Proof of Concept Demonstration Plan in the
\href{https://github.com/AHMEDELZARIA/squigglr/blob/main/pdfs/DevelopmentPlan/DevelopmentPlan.pdf}
{Development Plan} names four implementation risks. The demo addresses two
of them:

\begin{itemize}
    \item \textbf{Drawing evaluation.} One scoring method, read from each
    lesson's figure, must tell a complete drawing that meets the pass mark
    from an incomplete or inaccurate one, for every supported shape and at
    every drawing stage. If it cannot, no lesson can advance a child
    correctly.
    \item \textbf{Touchscreen and stylus interaction.} Input coordinates,
    stroke continuity and responsiveness may differ between devices and
    browsers.
\end{itemize}

The other two risks, editable drawing refinement and model-based creative
enhancement, belong to the creative space. The team deferred both until
after the POC, because the guided lessons do not depend on them and neither
experiment could be run properly in the time left. The supervisor agreed on
November 18 that this scope is enough for a proof of concept (issue \#370).
Editable refinement stays the main technical uncertainty after the demo.

\subsection*{What the team will demonstrate}

The demo, on November 25, shows the guided lessons running on an iPad, served
from a laptop on the same Wi-Fi. It covers experiments 1 (drawing input) and 4
(the square and circle lessons end to end) of the Development Plan:

\begin{itemize}
    \item An adult picks the Square or Circle lesson, sets the nine practice
    values and starts it.
    \item The child watches a demonstration, then draws through three stages
    with less guidance at each stage: trace the outline, endpoints only and
    no guide.
    \item Passing, failing and incomplete attempts show that each attempt is
    scored against the adult's values and that a pass moves the child to the
    next stage. A demo panel shows the result, the stage and the counters
    after each attempt.
    \item Drawing works with a stylus.
\end{itemize}

The square has only straight sides and the circle only curves, so the two
lessons together show that the scoring does not depend on the shape. The POC
has no accounts, child profiles, saved records, progress summary or creative
space, and is not hosted.

\section{Team Meeting Attendance}

\begin{table}[H]
\centering
\begin{tabular}{ll}
\toprule
\textbf{Student} & \textbf{Meetings}\\
\midrule
Total & 11\\
Rayan Nasrallah & 11\\
Ahmed Elzaria & 11\\
Muhammad Huzaifah & 10\\
Shamil Canbolat & 7\\
Moustafa Moustafa & 11\\
\bottomrule
\end{tabular}
\end{table}

The counts come from the 11 team meeting issues: \#6, \#7, \#8, \#9, \#11,
\#32, \#33, \#68, \#77, \#202 and \#323. The attendance for the meeting of
November 11 (\#323) was entered on November 19, when this report was
written. Every absence was announced in advance or had a reason that the
charter accepts.

\section{Supervisor/Stakeholder Meeting Attendance}

\noindent \textbf{Supervisor's Name: } Dr.~Christopher Anand

\begin{table}[H]
\centering
\begin{tabular}{ll}
\toprule
\textbf{Student} & \textbf{Meetings}\\
\midrule
Total & 5\\
Rayan Nasrallah & 5\\
Ahmed Elzaria & 4\\
Muhammad Huzaifah & 4\\
Shamil Canbolat & 4\\
Moustafa Moustafa & 4\\
\bottomrule
\end{tabular}
\end{table}

The counts come from the 5 supervisor meeting issues: \#21, \#107, \#201,
\#321 and \#370. All five members attended the first four meetings. For the
meeting of November 18 (\#370), the team agreed in advance that Rayan would
attend alone and present the POC. The team also interviewed three domain
specialists on October 1 (\#73). That issue has no attendance list, so the
table does not count it.

\section{Lecture Attendance}

\begin{table}[H]
\centering
\begin{tabular}{ll}
\toprule
\textbf{Student} & \textbf{Lectures}\\
\midrule
Total & 12\\
Rayan Nasrallah & 5\\
Ahmed Elzaria & 7\\
Muhammad Huzaifah & 3\\
Shamil Canbolat & 4\\
Moustafa Moustafa & 5\\
\bottomrule
\end{tabular}
\end{table}

The total is the instructor's figure of 12 lectures between the start of
classes and the POC. The counts come from the 8 lecture issues: \#1, \#2,
\#3, \#4, \#5, \#12, \#22 and \#76. Issue \#5 (September 17) records no
attendee. The team opened no lecture issue after October 1, so the table
counts every member as absent from the 4 lectures after that date.

\section{TA Document Discussion Attendance}

\noindent \textbf{TA's Name: } Christopher Schankula

\begin{table}[H]
\centering
\begin{tabular}{ll}
\toprule
\textbf{Student} & \textbf{Meetings}\\
\midrule
Total & 4\\
Rayan Nasrallah & 4\\
Ahmed Elzaria & 4\\
Muhammad Huzaifah & 4\\
Shamil Canbolat & 4\\
Moustafa Moustafa & 3\\
\bottomrule
\end{tabular}
\end{table}

The counts come from the 4 TA meeting issues: \#10, \#78, \#203 and \#322.
The one absence was announced in advance or had a reason that the charter
accepts.

\section{Commits, Lines Added and Lines Deleted, Pull Requests}

\begin{table}[H]
\centering
\makebox[\textwidth][c]{%
\begin{tabular}{lllllll}
\toprule
\textbf{Student} & \textbf{Commits} & \textbf{Percent} & \textbf{Lines Added} &
\textbf{Lines Deleted} & \textbf{Pull Requests}\\
\midrule
Total & 334 & 100\% & 54602 & 17126 & 146\\
Rayan Nasrallah & 90 & 26.9\% & 6114 & 2600 & 36\\
Ahmed Elzaria & 68 & 20.4\% & 21148 & 2343 & 33\\
Muhammad Huzaifah & 68 & 20.4\% & 10976 & 1503 & 34\\
Shamil Canbolat & 38 & 11.4\% & 10729 & 9043 & 16\\
Moustafa Moustafa & 70 & 21.0\% & 5635 & 1637 & 27\\
\bottomrule
\end{tabular}}
\end{table}

How the team counted:

\begin{itemize}
    \item Commits are the commits on \texttt{main} that are not merge
    commits. The table leaves out 202 merge commits and the 100 commits that
    the documentation workflow made to publish PDFs.
    \item Lines come from \texttt{git log --numstat} over the same commits.
    Binary files do not count.
    \item Pull requests are all pull requests that each member opened:
    125 merged, 10 open and 11 closed without a merge.
\end{itemize}

Two commits distort the line counts. Ahmed's first commit added the course
template (11354 lines added). One of Shamil's commits wrapped the SRS text at
80 columns (7088 lines added, 7953 deleted).

The POC application is not on \texttt{main} yet. It is in pull request
\#366 from the \texttt{POC} branch, which Rayan opened: 16 commits and
13332 lines added. Nine more pull requests with documentation fixes are
open: six from Shamil (6 commits) and three from Muhammad (6 commits).

No commit on \texttt{main} names two team members as authors. Eight commits
name an AI tool as co-author: five name Copilot Autofix and three name
Claude.

\section{Issue Tracker}

\begin{table}[H]
\centering
\begin{tabular}{lll}
\toprule
\textbf{Student} & \textbf{Authored (O+C)} & \textbf{Assigned (C only)}\\
\midrule
Rayan Nasrallah & 30 & 38 \\
Ahmed Elzaria & 30 & 45 \\
Muhammad Huzaifah & 57 & 36 \\
Shamil Canbolat & 0 & 13 \\
Moustafa Moustafa & 75 & 62 \\
\bottomrule
\end{tabular}
\end{table}

The repository has 228 issues. Team members authored 192 of them. Members of
two other teams authored the other 36 as peer reviews. The authored counts
include the meeting and lecture issues. An issue with several assignees
counts once for each assignee.

\section{CICD via GitHub Actions}

The team uses one GitHub Actions workflow today,
\href{https://github.com/AHMEDELZARIA/squigglr/blob/main/.github/workflows/latex-pages.yml}
{\texttt{latex-pages.yml}}, in the folder \texttt{.github/workflows}.

\begin{itemize}
    \item On each pull request into \texttt{main}, the workflow compiles the
    documents that the change affects. A LaTeX error fails the check before
    review.
    \item On each push to \texttt{main}, the workflow compiles the same
    documents, commits the PDFs to \texttt{pdfs/} and publishes them on
    \href{https://ahmedelzaria.github.io/squigglr/}{GitHub Pages}.
    \item Three scripts in
    \href{https://github.com/AHMEDELZARIA/squigglr/tree/main/.github/scripts}
    {\texttt{.github/scripts}} choose the documents to build and build the
    index page.
\end{itemize}

The workflow has run 444 times: 318 runs on pull requests, 125 on pushes and
1 manual run. 406 runs succeeded and 38 failed.

The POC application has 65 unit tests (Vitest), a type check and a
production build. The team runs them by hand for now. After the application
merges, the team will add the checks that the Development Plan lists:
formatting (Prettier), lint (ESLint), type checking (TypeScript), unit and
component tests with a coverage report (Vitest), a production build
(Next.js) and browser smoke tests (Playwright). GitHub Actions will run
these checks on pull requests. The team also plans preview deployments on
Vercel and a manual promotion to production.

\section{Custom Documents (CDs)}

\texttt{TeamComposition.csv} first listed a Literature Review Report for the
fall term and a Usability Report for the winter term.

The team replaced the fall document with a Specialist Interview Report. The
team interviewed three domain specialists on October 1, and the report
analyses which findings shaped which requirements and hazards. The
instructor approved the change through merge request 554 on the course
GitLab (issue \#164). The report was due with the Design Document Rev~$-1$
on November 16 and is in the repository
(\href{https://github.com/AHMEDELZARIA/squigglr/blob/main/pdfs/CDs/Interviews/InterviewReport.pdf}
{\texttt{InterviewReport.pdf}}).

The team keeps the Usability Report for the winter term and sees no reason
to change it. The POC gives the team a working lesson for the usability
sessions that the report plans.

\section{Team Charter Trigger Items}

\subsection*{Quantified triggers}

The charter is an appendix of the Development Plan. The team added it to the
repository on September 26. Any one of these four events triggers a workload
and accountability discussion:

\begin{enumerate}
    \item Two unexcused meeting absences within four weeks.
    \item Three arrivals more than ten minutes late within four weeks,
    without prior notice or an acceptable reason.
    \item Two missed internal deadlines within four weeks, without an agreed
    extension or an acceptable reason.
    \item No substantive progress on assigned work across two consecutive
    weekly reviews, without a communicated blocker or an agreed adjustment.
\end{enumerate}

The charter also sets these measurable expectations:

\begin{itemize}
    \item All five members attend every team, supervisor and teaching team
    meeting.
    \item Another member reviews each substantial contribution.
    \item Review-ready work is complete 48 hours before each course deadline.
    \item Each member contributes about ten hours a week.
    \item Fairness is judged on completed work, review participation,
    attendance and outstanding tasks. Commit counts alone do not decide it.
\end{itemize}

\subsection*{Violations}

The team's records show no violation of a trigger.

\begin{itemize}
    \item \textbf{Trigger 1.} Since September 26 the issues record 13
    meetings and 8 absences across the five members. Four of those absences
    are the supervisor meeting of November 18, which the team agreed that
    one member would attend. Every other absence was announced in advance
    or had a reason that the charter accepts.
    \item \textbf{Triggers 2 and 3.} The meeting issues do not record late
    arrivals, and the repository does not record internal deadlines apart
    from course deadlines. The team cannot count these two triggers from its
    records.
    \item \textbf{Trigger 4.} No issue records a workload and accountability
    discussion, which is the step that follows any trigger.
\end{itemize}

The metrics in this report also show three gaps against the charter and the
Development Plan. None of them is a trigger.

\begin{itemize}
    \item \textbf{Review participation is uneven.} Teammates reviewed 109 of
    the 125 merged pull requests (87\%), but one member did 51 of the
    reviews and another did 1. The other 16 pull requests merged without a
    teammate's review. The charter asks for a review of each substantial
    contribution.
    \item \textbf{The repository records fewer meetings than the Development
    Plan schedules.} The plan schedules three team meetings and one
    supervisor meeting each week. The repository records 11 team meetings
    and 5 supervisor meetings since the start of term.
    \item \textbf{The lecture record stops on October 1.} The charter sets
    no trigger for lectures, but this report needs the record.
\end{itemize}

\subsection*{Plan}

No violation needs a recovery plan. The team will close the gaps in its
records and in its reviews:

\begin{itemize}
    \item The notetaker will mark each absence in a meeting issue as excused
    or unexcused, and will note each arrival more than ten minutes late.
    The team can then count triggers 1 and 2 from the issues.
    \item Each task issue will state its internal deadline, as the charter
    already requires. The milestone QA coordinator will list missed internal
    deadlines in the milestone's tracking issue, so that the team can count
    trigger 3.
    \item The team will open a lecture issue for each remaining lecture and
    tick attendance on the same day.
    \item Authors will request reviews from the members with the fewest
    reviews first, and no pull request will merge without a teammate's
    review.
    \item At its next meeting, the team will decide whether to record every
    scheduled meeting or to change the meeting schedule in the Development
    Plan.
\end{itemize}

The team keeps the four triggers as they are. The triggers were not
ambiguous. The records were too thin to check them.

\section{Additional Productivity Metrics}

The team tracks code review as an additional metric, because the charter
names review participation as a measure of fair contribution.

\begin{table}[H]
\centering
\begin{tabular}{lll}
\toprule
\textbf{Student} & \textbf{PRs Merged} & \textbf{Teammates' PRs Reviewed}\\
\midrule
Total & 125 & 118\\
Rayan Nasrallah & 31 & 18\\
Ahmed Elzaria & 31 & 28\\
Muhammad Huzaifah & 30 & 20\\
Shamil Canbolat & 8 & 1\\
Moustafa Moustafa & 25 & 51\\
\bottomrule
\end{tabular}
\end{table}

PR means pull request. The review column counts each pull request once for
each teammate who submitted a review on it. Reviews from the Copilot review
tool do not count.

The team also tracks feedback issues, which carry the \texttt{review} label.
They hold the findings from peer reviews, the TA and the supervisor. Of 72
feedback issues, 40 are closed and 32 are open.

\end{document}
